\chapter{Sistemas Axiomáticos de Sheffer (1913)}

\section{Reducción de Operadores}

En 1913, Henry M. Sheffer demostró que el álgebra de Boole puede ser completamente definida utilizando un único operador lógico, en lugar de los tres (disyunción, conjunción y negación) requeridos por el sistema de Huntington. Los dos operadores capaces de esta universalidad funcional son el operador \textbf{NAND} (Barra de Sheffer, $\uparrow$) y el operador \textbf{NOR} (Flecha de Peirce, $\downarrow$). 

Todos los axiomas de Huntington (1904) de los que hemos partido son reducibles a un grupo más pequeño y estricto de axiomas basados exclusivamente en uno de estos dos operadores.

\section{Axiomatización mediante NAND (Barra de Sheffer)}

El sistema axiomático propuesto por Sheffer para la operación NAND consta de los siguientes cinco postulados sobre una clase $K$:

\begin{teorema}{Postulados de Sheffer (1913) - Operador NAND}{sheffer_nand}
\begin{enumerate}
    \item Existen al menos dos elementos distintos en $K$.
    \item Clausura: Para cualesquiera $a, b \in K$, el resultado de $a \uparrow b$ también pertenece a $K$.
    \item $(a \uparrow a) \uparrow (a \uparrow a) = a$
    \item $a \uparrow (b \uparrow (b \uparrow b)) = a \uparrow a$
    \item $(a \uparrow (b \uparrow c)) \uparrow (a \uparrow (b \uparrow c)) = ((b \uparrow b) \uparrow a) \uparrow ((c \uparrow c) \uparrow a)$
\end{enumerate}
\end{teorema}

Es directo demostrar que los axiomas de Huntington implican estos cinco postulados:

\begin{proof}
Los postulados 1 y 2 son inmediatos y se asumen de base, garantizando la existencia de los elementos y la definición del operador binario.

Para demostrar el postulado 3, usamos la equivalencia $x \uparrow x = \neg x$:
\begin{align*}
(a \uparrow a) \uparrow (a \uparrow a) &= \neg a \uparrow \neg a & \text{Definición de } \uparrow \\
&= \neg (\neg a) & \text{Definición de } \uparrow \\
&= a & \text{Involución (Doble Negación)}
\end{align*}

Para el postulado 4, partimos del lado izquierdo sabiendo que $x \uparrow 1 = \neg x$:
\begin{align*}
a \uparrow (b \uparrow (b \uparrow b)) &= a \uparrow (b \uparrow \neg b) & \text{Definición de } \uparrow \\
&= a \uparrow \neg (b \wedge \neg b) & \text{Definición de } \uparrow \\
&= a \uparrow \neg (\bot) & \text{Axioma de Complementarios ($Comp$)} \\
&= a \uparrow \top & \text{Complemento del Neutro} \\
&= \neg (a \wedge \top) & \text{Definición de } \uparrow \\
&= \neg a & \text{Axioma de Elemento Neutro ($E_n$)} \\
&= a \uparrow a & \text{Definición de } \uparrow
\end{align*}

Para el postulado 5, desarrollamos ambas partes de la igualdad hasta alcanzar una expresión booleana idéntica. Empezamos por el lado izquierdo:
\begin{align*}
(a \uparrow (b \uparrow c)) \uparrow (a \uparrow (b \uparrow c)) &= \neg (a \uparrow (b \uparrow c)) & \text{Por postulado 3} \\
&= \neg (\neg (a \wedge (b \uparrow c))) & \text{Definición de } \uparrow \\
&= a \wedge (b \uparrow c) & \text{Involución} \\
&= a \wedge \neg (b \wedge c) & \text{Definición de } \uparrow \\
&= a \wedge (\neg b \vee \neg c) & \text{Leyes de De Morgan}
\end{align*}

Desarrollamos ahora el lado derecho:
\begin{align*}
((b \uparrow b) \uparrow a) \uparrow ((c \uparrow c) \uparrow a) &= (\neg b \uparrow a) \uparrow (\neg c \uparrow a) & \text{Definición de } \uparrow \\
&= \neg (\neg b \wedge a) \uparrow \neg (\neg c \wedge a) & \text{Definición de } \uparrow \\
&= (b \vee \neg a) \uparrow (c \vee \neg a) & \text{De Morgan e Involución} \\
&= \neg ((b \vee \neg a) \wedge (c \vee \neg a)) & \text{Definición de } \uparrow \\
&= \neg (\neg a \vee (b \wedge c)) & \text{Axioma de Distributividad ($Dist_\vee$)} \\
&= a \wedge \neg (b \wedge c) & \text{De Morgan e Involución} \\
&= a \wedge (\neg b \vee \neg c) & \text{Leyes de De Morgan}
\end{align*}
Al llegar ambos desarrollos a la misma expresión booleana ($a \wedge (\neg b \vee \neg c)$), la igualdad queda estrictamente demostrada.
\end{proof}

\section{Axiomatización mediante NOR (Flecha de Peirce)}

El razonamiento es rigurosamente simétrico si se escoge el operador dual $\downarrow$.

\begin{teorema}{Postulados de Sheffer Duales - Operador NOR}{sheffer_nor}
\begin{enumerate}
    \item Existen al menos dos elementos distintos en $K$.
    \item Clausura: Para cualesquiera $a, b \in K$, el resultado de $a \downarrow b$ pertenece a $K$.
    \item $(a \downarrow a) \downarrow (a \downarrow a) = a$
    \item $a \downarrow (b \downarrow (b \downarrow b)) = a \downarrow a$
    \item $(a \downarrow (b \downarrow c)) \downarrow (a \downarrow (b \downarrow c)) = ((b \downarrow b) \downarrow a) \downarrow ((c \downarrow c) \downarrow a)$
\end{enumerate}
\end{teorema}

\begin{proof}
La demostración sigue un patrón idéntico de dualidad. El postulado 3 se demuestra recordando que $x \downarrow x = \neg x$:
\begin{align*}
(a \downarrow a) \downarrow (a \downarrow a) &= \neg a \downarrow \neg a = \neg (\neg a) = a
\end{align*}

En el postulado 4, sabiendo que $x \downarrow 0 = \neg x$, el término interior $b \downarrow (b \downarrow b)$ se evalúa a $\bot$, y por tanto:
\begin{align*}
a \downarrow (b \downarrow (b \downarrow b)) &= a \downarrow (b \downarrow \neg b) = a \downarrow \neg (b \vee \neg b) = a \downarrow \neg (\top) = a \downarrow \bot \\
&= \neg (a \vee \bot) = \neg a = a \downarrow a
\end{align*}

Para el postulado 5, el proceso dual lleva ambos lados de la ecuación a la forma idéntica $a \vee (\neg b \wedge \neg c)$.
\end{proof}
